% =====================================================================================
% DRAFT: smoother transition into Section 2.5 "Programmable photonic systems"
% (not included in main.tex; ch2-theory.tex is not modified)
%
%   [1] Intro of Section 2.5                 -> REPLACES the current 3-sentence intro
%   [2] First sentence of 2.5.1 (optional)   -> REPLACES "The end-to-end insertion loss of an optical path is ..."
%   [3] Last paragraph of 2.4 (optional)     -> REPLACES "Building on these developments, we have proposed ..."
%   [4] Chapter 2 introduction (full text)   -> REPLACES the whole chapter intro (lines 4-15 of ch2-theory.tex)
% =====================================================================================


% -------------------------------------------------------------------------------------
% [1] Intro of Section 2.5  (recommended)
%
% Idea of the transition: Section 2.4 ends with the SDN control of optical devices.
% An SDN controller can only act on the resources that a device exposes, and for a
% programmable photonic device these resources go beyond the PUCs of the photonic core.
% This leads naturally to the system-level components and to the two platforms.
% -------------------------------------------------------------------------------------

\section{Programmable photonic systems}\label{sub:pip_systems}
The SDN framework described above controls the physical layer through the resources that each device exposes to the controller.
For a programmable photonic device, these resources are not limited to the PUCs of the photonic core described in Sections~\ref{sec:functional_blocks} and~\ref{sec:mesh_architectures}.
In a deployable system, the photonic core is complemented by system-level components that couple the light into and out of the chip, handle its polarization, compensate its insertion loss, and monitor its optical power, together with the electronics and software that drive them.
These components often dominate the end-to-end performance of the system, and they also determine which parameters can be configured and observed by the control plane.
This section first summarizes these components and then presents the two iPronics platforms used in this thesis: Smartlight, used in Chapter~\ref{chap:applications_on_pip}, and iPronics ONE, used in Chapter~\ref{chap:sdn_for_ocs}.


% -------------------------------------------------------------------------------------
% [2] First sentence of Subsection 2.5.1 (optional)
%
% The subsection currently starts directly with the loss equation. This sentence
% introduces it from the network point of view. Keep the rest of the paragraph,
% starting from the equation, unchanged.
% -------------------------------------------------------------------------------------

\subsection{System-level components}\label{sub:system_component}
From the network point of view, the first figure of merit of a programmable photonic system is the end-to-end insertion loss of an optical path.
This loss is the sum of a fixed contribution, given by the fiber-to-chip couplers, the polarization-handling components, and the waveguide propagation, and of a contribution that grows with the size of the circuit:
% ... equation eq:ch2-loss-budget and the rest of the paragraph unchanged ...


% -------------------------------------------------------------------------------------
% [3] Last paragraph of Section 2.4 "Optical SDN" (optional)
%
% Present tense and "this thesis" instead of "we have proposed", so that the section
% hands over to 2.5 instead of jumping directly to Chapter 5.
% -------------------------------------------------------------------------------------

Building on these developments, this thesis proposes an SDN control layer for silicon photonic OCS, whose design, implementation, and evaluation are presented in Chapter~\ref{chap:sdn_for_ocs}.
Such a control layer relies on the hardware resources that the OCS exposes, which are introduced in the next section.


% -------------------------------------------------------------------------------------
% [4] Chapter 2 introduction (full text)
%
% Replace everything between "\chapter{...}\label{chap:fundamentals}" and
% "\section{Basic building blocks}" with the text below.
%
% Changes with respect to the current intro:
%  - one paragraph per section, in the actual order 2.1 -> 2.5, each with \ref to its section;
%  - building blocks listed in the order of Section 2.1 (splitter, phase shifter, PUC, crossing);
%  - Section 2.5 (programmable photonic systems) is now announced;
%  - each part says where it is used later (Chapters 3, 4, 5), which helps the reader
%    follow the later developments;
%  - hyphenation: device-level, large-scale, graph-theoretic, higher-layer, model-driven.
% -------------------------------------------------------------------------------------

The rapid advancement of silicon photonics has transitioned integrated optical systems from static designs toward reconfigurable, software-defined architectures.
This evolution enables the dynamic manipulation of optical signals, allowing a single hardware platform to support reconfigurable networking functions, adaptive interconnects, and scalable optical switching \cite{bogaerts2020programmable}.
As programmable integrated photonic circuits continue to increase in scale and complexity, efficient methodologies for circuit abstraction, automated configuration, and software-defined control become essential for practical deployment.
This chapter introduces these foundations following the bottom-up order of the cross-layer structure of Figure~\ref{fig:ch1-thesis}: from the photonic devices, through their abstraction and programming, up to their control within a network.

Section~\ref{sec:functional_blocks} begins with the fundamental building blocks, including 50:50 power splitters, phase shifters, 2$\times$2 programmable unit cells (PUCs), and waveguide crossings, and presents the transfer-matrix models that describe their optical behavior.
Building upon these device-level elements, Section~\ref{sec:mesh_architectures} introduces the programmable mesh topologies constructed from interconnected unit cells, namely feedforward and recirculating meshes.

To address the growing complexity of large-scale programmable photonic fabrics, Section~\ref{sec:graph} develops a graph-theoretic abstraction that represents photonic circuits as weighted graphs composed of nodes and edges, and describes how optical paths are computed on these graphs with shortest-path algorithms.
This abstraction is the basis of the automated routing methods of Chapter~\ref{chap:applications_on_pip} and of the architecture analysis of Chapter~\ref{chap:ocs_archs}.

In parallel with hardware abstraction and programming, the scalable operation of programmable photonic systems also requires integration with higher-layer network control frameworks.
Section~\ref{sec:sdn} therefore introduces the principles of software-defined networking (SDN), model-driven management with YANG, and southbound communication interfaces such as gNMI, which form the foundation of the control framework developed in Chapter~\ref{chap:sdn_for_ocs}.

Finally, Section~\ref{sub:pip_systems} describes the system-level components that turn a programmable photonic chip into a deployable system, such as optical amplification, polarization handling, and optical power monitoring, and presents the two platforms on which the experiments of this thesis are carried out: Smartlight and iPronics ONE.
